\documentclass{article}
\usepackage{mathnotes}

\title{Real Numbers}
\description{Axiomatically defines the real numbers as an ordered complete field, covering field axioms, order properties, completeness, supremum and infimum, and the absolute value function.}

\begin{document}

\section{Real Numbers}

\begin{definition}[Real Numbers]\synonyms{reals}
The real numbers $\mathbb{R}$ are a set of objects along with two binary operations $+ : \mathbb{R} \times \mathbb{R} \to \mathbb{R}$ and $\cdot : \mathbb{R} \times \mathbb{R} \to \mathbb{R}$ that satisfy the 9 field axioms along with the Order axiom and the Completeness axiom. That is, the reals are an ordered, complete field.
\end{definition}

We know that a \pagelink[field is a commutative division ring]{algebra/abstract/rings}.

The \textbf{Order axiom} states that there exists a subset of $\mathbb{R}^+$ of $\mathbb{R}$ such that for all $a, b \in \mathbb{R}^+,a + b \in \mathbb{R}^+, a \cdot b \in \mathbb{R}^+.$ and for all $a \in \mathbb{R}$ exactly one of the following is true: (i) $a \in \mathbb{R}^+$, (ii) $-a \in \mathbb{R}^+$, or (iii) $a = 0$.

A more familiar and equivalent way of stating this can be achieved by defining $a < b$ to mean $b + (-a) \in \mathbb{R}^+$, and $a > b$ to mean $a + (-b) \in \mathbb{R}^+$, for $a, b, c \in \mathbb{R}$,

\begin{itemize}
\item Exactly one of the following is true: $a < b, ~ a > b,$ or $a = b$.
\end{itemize}

\begin{itemize}
\item If $a < b$ and $c > 0$, we have $a \cdot c < b \cdot c.$
\end{itemize}

\begin{itemize}
\item If $a < b$ and $b < c$ then $a < c$ (transitivity).
\end{itemize}

Note that rational numbers also satisfy the field axioms and the Order axiom. Complex numbers satisfy the field axioms but not the Order axiom because there is no total ordering that can be defined on complex numbers.

To give the Completeness axiom, we must first define what a least upper bound is.

\begin{definition}[Upper bound]
The number $m$ is said to be an \textbf{upper bound} of a nonempty \dref{set} $A$ if $x \leq m$ for all $x \in A$.
\end{definition}

\begin{definition}[Least upper bound]\synonyms{supremum}
The number $L$ is said to be the \textbf{least upper bound} of the set $A$ if it's an \dref{upper-bound} of $A$ and if $L \leq m$ for all \dref{upper-bounds} $m$ of $A$.
 
The least upper bound is also known as the \textbf{supremum.}
\end{definition}

\begin{note}
Neither an \dref{upper-bound} nor a \dref{least-upper-bound} need to be in $A$.
\end{note}

\begin{definition}[Lower bound]
The number $m$ is said to be a \textbf{lower bound} of a nonempty \dref{set} $A$ if $x \geq m$ for all $x \in A.$
\end{definition}

\begin{definition}[Greatest lower bound]\synonyms{infimum}
There is also a concept of a \textbf{greatest lower bound}, which is a \dref{lower-bound} that is greater than or equal to every other \dref{lower-bound}. The greatest lower bound is also known as the \textbf{infimum.}
\end{definition}

The \textbf{Completeness axiom} states that every nonempty subset $A$ of $\mathbb{R}$ that's bounded above has a least upper bound.

One consequence of the Completeness axiom is that the real numbers contain numbers that are not rational. For example, the subset of $\mathbb{R}$ that is the finite decimal approximations of $\sqrt{5}$ is,

\[ \{2, 2.2, 2.23, 2.236, \dots \}, \]

must have a supremum due to the Completeness axiom. That supremum is $\sqrt{5}$, but $\sqrt{5}$ is not rational, therefore, the reals contain non-rational numbers.

The \textbf{absolute value} of a real number is defined as follows:

\[ |x| = \begin{cases} x & \text{if } x \geq 0 \\ -x & \text{if } x < 0 \end{cases}. \]

One fact that follow from this definition is that for $a, b \in \mathbb{R}$, if $\|a\| < b$, then $-b < a < b$. \emph{Proof}: We have two cases to consider:

\begin{itemize}
\item If $a >= 0$, then $\|a\| = a$, so $a < b$. Since $a$ is positive, $b$ must also be positive, and therefore $-b$ must be negative, and so $-b < a < b$.
\end{itemize}

\begin{itemize}
\item If $a < 0$, then $\|a\| = -a$, so $-a < b$. Since $a$ is negative, $-a$ must be positive, so $b$ must also be positive, and therefore $a < b$. Now, $-a < b \iff a > -b \iff -b < a$, and therefore $-b < a < b.$ $\square$
\end{itemize}

Another fact is that the absolute value function is a norm, and thus satisfies the triangle inequality. For real $a, b$:

\[ |a + b| \leq |a| + |b| \]

\begin{definition}[Archimedean]
The Archimedean property is that given two positive numbers $x$ and $y,$ there is an integer $n$ such that $nx > y.$
\end{definition}

\begin{theorem}\label{reals-are-archimedean}
The real numbers are Archimedean.
\end{theorem}

\begin{theorem}\label{rationals-are-dense-in-reals}
The rationals are \dref{dense} in the reals.
\end{theorem}

\begin{proof}
Let $x \in \mathbb{R}$ and $\epsilon > 0.$ Pick $n \in \mathbb{N}$ such that $1/n < \epsilon.$ Now, because \dref[the reals are Archimedean]{reals-are-archimedean}, for some $m \in \mathbb{Z}$ we have that $m = \lfloor nx \rfloor.$ Then,

\[ m \leq nx \leq m+1, \quad \frac{m}{n} \leq x \leq \frac{m+1}{n}. \]

Now, if we let $q = \frac{m}{n},$ we have

\[ 0 \leq x - q < \frac{1}{n} < \epsilon, \]

so $|x - q| < \epsilon.$ Therefore, every \dref{neighborhood} of $x$ contains some $q \in \mathbb{Q},$ and so $x$ is either a limit point of $\mathbb{Q}$ or else $x \in \mathbb{Q},$ and so the rationals are therefore dense in the reals.
\end{proof}

\begin{corollary}\label{rational-between-any-two-reals}
If $x, y \in \mathbb{R},$ and $x < y,$ then there exists a $p \in \mathbb{Q}$ such that $x < p < y.$ That is, there is always a rational number between any two distinct real numbers.
\end{corollary}

\begin{proof}
Pick $n \in \mathbb{N}$ such that $n(y - a) > 1.$ Because because \dref[the reals are Archimedean]{reals-are-archimedean}, there is some $k \in \mathbb{Z}$ such that $nx < k < ny$ and $x < n/k < y.$ Thus, $q = n/k \in \mathbb{Q}$ and $r \in (x, y).$
\end{proof}

\end{document}
